% Source: https://artofproblemsolving.com/community/c3h2360552_objectzs_1000th_post_problems_and_other_special_posts
% Stable identity: aops-topic-2360552-q09
The mode is $2$ and the mean is $\dfrac{25+x}{7}$.
\begin{itemize}
\item $x\le2$: the median is $2$, equal to the mode. A non-constant arithmetic progression cannot repeat a term, so this is impossible.
\item $2<x<4$: the median is $x$ and the mean lies between $\dfrac{27}{7}$ and $\dfrac{29}{7}$. The order $2,x,\dfrac{25+x}{7}$ needs $2x=2+\dfrac{25+x}{7}$, so $x=3$ (giving $2,3,4$). The order $2,\dfrac{25+x}{7},x$ needs $x=\dfrac{36}{5}$, which is not below $4$.
\item $x\ge4$: the median is $4$ and the mean exceeds $4$, so the progression is $2,4,6$ and $\dfrac{25+x}{7}=6$ gives $x=17$.
\end{itemize}
The sum is $3+17=\boxed{20}$.
